% from ega4-18-fr.tex
\oldpage[IV-4]{143}

\begin{proposition}[18.7.2]
\label{IV.18.7.2}
Soit $A$ un anneau local noethérien. Pour que ${}^h\!A$ soit un $\mathbf P$-anneau (7.3.13), il faut et il suffit que $A$ le soit.
\end{proposition}

En effet, en vertu de (18.6.6, (iv)), le complété $\widehat A$ de $A$ est aussi le complété de ${}^h\!A$ et l'on a donc $A\subset{}^h\!A\subset\widehat A$. D'après (18.6.9), pour tout $x\in\Spec(A)$, la fibre en $x$ du morphisme $\Spec({}^h\!A)\to\Spec(A)$ est discrète et finie, et en chacun des points $y_i$ de cette fibre, $k(y_i)$ est une extension algébrique séparable de $k(x)$. La fibre formelle de $A$ au point $x$ est donc un préschéma \emph{somme} des fibres formelles de ${}^h\!A$ aux points $y_i$. La conclusion résulte alors de l'hypothèse faite sur $\mathbf P$ dans (18.7.1).

\begin{corollary}[18.7.3]
\label{IV.18.7.3}
Soit $A$ un anneau local noethérien. Pour que ${}^h\!A$ soit universellement japonais, il faut et il suffit que $A$ le soit.
\end{corollary}

En effet, pour un anneau local noethérien $B$, il revient au même de dire que $B$ est universellement japonais, ou que les fibres formelles de $B$ sont géométriquement réduites (7.6.4 et 7.7.1); la conclusion résulte donc de (18.7.2).

\begin{corollary}[18.7.4]
\label{IV.18.7.4}
Soit $A$ un anneau local noethérien. Pour que les fibres formelles de ${}^h\!A$ soient géométriquement régulières, il faut et il suffit que celles de $A$ le soient.
\end{corollary}

C'est un cas particulier de (18.7.2).

\begin{proposition}[18.7.5]
\label{IV.18.7.5}
Soit $A$ un anneau local noethérien. Si $A$ est universellement caténaire (5.6.2) (resp. formellement caténaire (7.1.9), resp. strictement formellement caténaire (7.2.6)), il en est de même de ${}^h\!A$.
\end{proposition}

Avec les notations de (18.6.4), si $A$ est universellement caténaire, il en est de même des $A_\alpha$, qui sont des $A$-algèbres essentiellement de type fini (5.6.3). Pour prouver que dans ce cas ${}^h\!A$ est universellement caténaire, il suffira donc d'établir le lemme suivant:

\begin{lemma}[18.7.5.1]
\label{IV.18.7.5.1}
Soit $(B_\alpha,\varphi_{\beta\alpha})$ un système inductif d'anneaux noethériens et supposons que, pour $\alpha\leq\beta$, $\varphi_{\beta\alpha}$ fasse de $B_\beta$ un $B_\alpha$-module fidèlement plat et que $B=\varinjlim B_\alpha$ soit noethérien. Alors, si les $B_\alpha$ sont caténaires (resp. universellement caténaires), il en est de même de $B$.
\end{lemma}

Dire que $B$ est universellement caténaire équivaut à dire que $B[T_1,\ldots,T_n]$ est caténaire pour tout $n$ (5.6.2), et il est immédiat que le système inductif $(B_\alpha[T_1,\ldots,T_n])$ vérifie les mêmes conditions que $(B_\alpha)$; il suffit donc de prouver que si les $B_\alpha$ sont caténaires, il en est de même de $B$. Si $\mathfrak p_0\supset\mathfrak p_1\supset\cdots\supset\mathfrak p_r$ est une chaîne saturée d'idéaux premiers de $B$, les images réciproques $\mathfrak p_{0\alpha}\supset\mathfrak p_{1\alpha}\supset\cdots\supset\mathfrak p_{r\alpha}$ de ces idéaux dans $B_\alpha$ forment une chaîne d'idéaux premiers, et il suffit de montrer que pour $\alpha$ assez grand, cette chaîne est saturée. Or, chaque $\mathfrak p_i$ est un idéal de type fini, donc il existe un $\lambda$ tel que, pour $\alpha\geq\lambda$, on ait $\mathfrak p_i=\mathfrak p_{i\alpha}.B$. On a donc
\[
1=\operatorname{codim}(V(\mathfrak p_i),V(\mathfrak p_{i+1}))
=\operatorname{codim}(V(\mathfrak p_{i\alpha}),V(\mathfrak p_{i+1,\alpha}))
\]
puisque $B$ est un $B_\alpha$-module plat ((6.1.4) et ($0_{\mathrm I}$, 6.2.3)), ce qui achève de prouver le lemme.

Supposons maintenant $A$ formellement caténaire. Soient $\mathfrak p$ un idéal premier de ${}^h\!A$, $\mathfrak q$ son image réciproque dans $A$, de sorte que $\mathfrak q.{}^h\!A\subset\mathfrak p$, et par suite ${}^h\!A/\mathfrak p$ est un anneau quotient de ${}^h\!A/\mathfrak q.{}^h\!A$. Il suffit donc de prouver que ${}^h\!A/\mathfrak q.{}^h\!A$ est formellement caténaire pour tout idéal premier $\mathfrak q$ de $A$ (7.1.9), et comme ${}^h\!A/\mathfrak q.{}^h\!A={}^h(A/\mathfrak q)$ par (18.6.8), il suffit (7.1.11) de voir que si $A$ est \emph{intègre et formellement caténaire}, ${}^h\!A$ est \emph{formellement équidimensionnel}. Mais comme le complété de ${}^h\!A$ est égal à $\widehat A$, cela résulte de l'hypothèse que $A$ est formellement caténaire.

